\section{Discussion}
\label{sec:discussion}
% Written from the final tables (results/final/stats_held.json, results/bdcxr/stats.json). Failed hypotheses are named first.
\paragraph{Principal findings.} Three pre-specified results did not hold: the soft-tissue image alone did not improve any downstream task
(Table~\ref{tab:downstream_cls}, Table~\ref{tab:downstream_det}), SSIM fell short of the reference model by more than its margin
(Table~\ref{tab:h2}), and conformal intervals calibrated on synthetic data covered 14\,\% instead of 90\,\% of real bone pixels
(Table~\ref{tab:reliability}); recalibrated on 24 real pairs in an exploratory analysis, they covered 89\,\% (Table~\ref{tab:bone_recal}). Real-nodule contrast was mixed: higher than with the reference model on JSRT, lower on Node21
(Table~\ref{tab:nodule_cnr}). The central results held. Lesion absorption fell steadily with the twin weight (Table~\ref{tab:dose_response}); the
twin ensemble absorbed a third of what the released CT-supervised model absorbed and less than a third of what the matched control absorbed, on
an audit set no model had seen during development (Table~\ref{tab:audit_contrasts}); and it did so at non-inferior bone-region fidelity. The effect
replicated with two other encoders (Table~\ref{tab:encoders}) and on paediatric radiographs from a population no model was trained on (Section~\ref{sec:bdcxr}).

\paragraph{What the audit adds.} Fidelity metrics and lesion absorption rank suppressors differently (Figure~\ref{fig:pareto}). The v1 model has
the highest bone-region PSNR of the CT-supervised models and the highest absorption; ours and the reference model are within 0.06\,dB of each
other and a factor of three apart in absorption. BS-LDM, trained on dual-energy pairs from the same hospital, is about 7\,dB ahead in fidelity on
that hospital's images, yet in the development audit it absorbed 0.42 of an inserted lesion (at site positions not paired with the other
methods, so only descriptive). A suppressor that matches the dual-energy target closely can still remove lesions: fidelity to that target does
not protect what the target never showed it. Absorption must nevertheless be read beside suppression strength. The reference model removes more
rib-edge energy than ours (48\,\% against 30\,\%) and more non-rib energy (22\,\% against 8\,\%), so part of its higher absorption may reflect
stronger suppression in general; ours absorbs a third as much while removing about 60\,\% of the rib-edge energy.

\paragraph{What the regulariser costs.} On real paired radiographs the twin costs no measurable bone-region fidelity up to $w=16$ (21.22 to
21.28\,dB across $w=0\ldots16$); at $w=32$ and 64 the PSNR drops by about 0.3\,dB, so $w=16$ is the usable
operating point, chosen on development data before the freeze. The cost appears where the truth is exact: on the Qure-projection test split the
bone error of the twin ensemble is 41\,\% higher than that of the control (MAE 0.0041 against 0.0029), and in the ablation the validation error rises
from 0.0033 to 0.0057. The regulariser therefore trades some synthetic bone accuracy for lesion invariance; on real images that trade was not
visible. The ablation shows what carries the effect: gradients through both twin branches and a consistency term confined to the pasted lesion.
Showing lesions during training without the twin term (CLIA) did not lower absorption.

\paragraph{Relation to other methods.} The reference model is a SwinV2-B at $1024$~px trained with a far larger sample budget; we compare its
released weights as a user would run them, and we cannot say how much of the absorption difference is due to our regulariser rather than to
training differences, except through the matched no-twin control, which absorbs slightly more than the reference model (0.178 against 0.149). Diffusion models trained on real
dual-energy pairs \citep{sun2026bsldm,sun2026debonedit} answer a different question: they need paired data from the target site, sample slowly
(about 2.5 minutes per image for BS-LDM on a T4 GPU) and were not designed for lesion invariance.

\paragraph{Advice for evaluation and use.} We suggest that a bone suppressor be audited for lesion absorption before it is used, with the audit
reported next to a suppression measure; that suppressed images be offered beside the original rather than instead of it, which is also the only
use in which they helped downstream models here; and that the ensemble spread be used as a warning flag rather than as a calibrated interval: it
rose 2.7-fold on paediatric radiographs, but intervals calibrated on synthetic data did not cover real errors, and recalibration needed real
pairs. Whether classifiers that combine the radiograph with the soft-tissue and bone images in separate streams gain from suppression is tested
in \clsstudy.

\paragraph{Future work.} Prediction at $1024$~px, combining CT-derived and real dual-energy supervision with the twin term, conformal
recalibration on real pairs from more sites, lesion types other than nodules (consolidation, as in the paediatric pneumonia set), paediatric training data, and a
reader study that measures detectability rather than signal.
